\documentclass[
  spanish,
  letterpaper, oneside
]{article}

\def\documenttitle {Diversidad cultural y lingüística en el derecho}
\def\documentsubtitle {Interpretación intercultural y acceso a la justicia}
\def\documentsubject {Actividad 2 - Interaprendizaje en ambientes virtuales}

\def\documentauthor {Martin Jonathan de la Cruz}
\def\coursename {Interaprendizaje en ambientes virtuales}
\def\coursecode {004}

\def\universityname {Universidad Abierta y a Distancia de México}
\def\universityfaculty {Licenciatura en Derecho}
\def\universitydepartment {Interaprendizaje en ambientes virtuales}
\def\universitydepartmentimage {departamentos/UnADM}
\def\universitydepartmentimagecfg {height=1.57cm}
\def\universitylocation {Roma Norte, Ciudad de México}

\def\authortable {
  \begin{tabular}{ll}
    \textbf{Alumno:}
    & \begin{tabular}[t]{l}
      Martin Jonathan de la Cruz \\
    \end{tabular} \\
    Matrícula: & ES2611202040 \\
    Figura docente: & Janeth Santana Ramírez \\
    Semestre/Bloque: & 2 / 1 \\
    & \\
    \multicolumn{2}{l}{Fecha de realización: \today} \\
    \multicolumn{2}{l}{\universitylocation}
  \end{tabular}
}

\input{template}
\usepackage{helvet}
\renewcommand{\familydefault}{\sfdefault}
\usepackage{array}
\usepackage{longtable}
\usepackage{pdflscape}
\usepackage{tikz}
\usetikzlibrary{arrows.meta,positioning,calc,matrix,fit,shapes.geometric,shadows.blur}
\setcitestyle{authoryear,open={(},close={)}}

\def\coverwatermarkenabled {true}
\def\coverwatermarkimage {img/departamentos/UnADM.pdf}
\def\coverwatermarkopacity {0.16}
\def\coverwatermarkwidth {0.72\paperwidth}
\def\coverwatermarkxshift {0cm}
\def\coverwatermarkyshift {-0.35cm}

\newcommand{\insertcoverwatermark}{%
  \ifthenelse{\equal{\coverwatermarkenabled}{true}}{%
    \AddToShipoutPictureBG*{%
      \begin{tikzpicture}[remember picture,overlay]
        \node[
          opacity=\coverwatermarkopacity,
          inner sep=0pt,
          xshift=\coverwatermarkxshift,
          yshift=\coverwatermarkyshift
        ] at (current page.center) {%
          \includegraphics[width=\coverwatermarkwidth]{\coverwatermarkimage}%
        };
      \end{tikzpicture}%
    }%
  }{}%
}

\definecolor{unadmGreenDark}{HTML}{174A3A}
\definecolor{unadmGreen}{HTML}{5F8F3A}
\definecolor{unadmGold}{HTML}{B88A2A}
\definecolor{unadmPaper}{HTML}{F6F7F2}
\definecolor{unadmInk}{HTML}{1F2A24}

% -----------------------------------------------------------------------------
% AulaTeX:foro-producto v1 -- bloque de participación del foro (producto)
% Caja institucional con texto JUSTIFICADO, numeración a la izquierda (enumitem),
% contenido copiable y un botón "Copiar participación" que adjunta un .txt.
% -----------------------------------------------------------------------------
\usepackage[most]{tcolorbox}
\usepackage{attachfile}
\usepackage{enumitem}

% attachfile: el 'color' debe darse como TRIPLE RGB numérico (0..1); un nombre
% \definecolor[HTML] provoca 'Argument of \atfi@textcolor has an extra }'.
\attachfilesetup{color={0.373 0.561 0.227},print=false}

% Botón de copia: appearance={} oculta el PushPin y muestra el 2.o argumento como
% botón visible. Se evita fontawesome (ps2pk/gsftopk no disponibles en pdflatex).
\newcommand{\foroCopyButton}[1]{%
  \textattachfile[appearance={},description={Participacion del foro (texto copiable)},mimetype=text/plain]{#1}{%
    \colorbox{unadmGreen}{\small\textcolor{white}{\,Copiar participación\,}}%
  }%
}

% Entorno visual del foro (banda verde, fondo claro, texto justificado)
\newtcolorbox{forobox}[1][]{%
  enhanced, breakable,
  colback=unadmPaper, colframe=unadmGreen,
  boxrule=0pt, leftrule=3pt, arc=1.2mm,
  left=4mm, right=4mm, top=3mm, bottom=3mm,
  fonttitle=\bfseries\color{white}, coltitle=white, colbacktitle=unadmGreenDark,
  attach boxed title to top left={xshift=4mm,yshift=-2mm},
  boxed title style={colframe=unadmGreenDark,arc=0.6mm},
  #1
}

\begin{document}

\insertcoverwatermark
\templatePortrait
\templatePagecfg
\onehalfspacing

\begin{abstractd}
  La diversidad cultural y lingüística funciona como criterio de interpretación
  jurídica: los estándares de lo ``válido'', lo ``confiable'' y lo ``bien
  explicado'' están culturalmente situados y pueden hacer que una narración
  comunitaria u oral parezca insuficiente. A partir de un caso en que una
  estudiante de Derecho descarta un relato por no ajustarse al formato documental,
  se argumenta que la competencia intercultural no debilita el rigor probatorio,
  sino que exige investigar, contextualizar y traducir antes de concluir. El
  análisis se sostiene en la comunicación intercultural, el protocolo de
  perspectiva intercultural y el marco constitucional y de derechos lingüísticos, y
  se materializa en una participación argumentada y una retroalimentación entre
  pares.
\end{abstractd}

\templateIndex
\templateFinalcfg

\section{Introducción}

La relación entre cultura, diversidad cultural, diversidad lingüística y formación
profesional en Derecho parte de una premisa: todas las personas interpretamos la
realidad desde referentes culturales y lingüísticos. Por ello, las formas de
hablar, comprender, explicar y valorar una situación influyen en cómo se interpreta
un problema jurídico. El propósito no es sustituir el rigor por sensibilidad
genérica, sino reconocer que la comprensión de los hechos también es una operación
cultural que la persona profesional del Derecho debe examinar de manera crítica.

El caso detonante presenta a una estudiante que atiende a una persona con un
conflicto por un terreno familiar. El acuerdo se hizo de palabra ante la comunidad,
con participación de la familia y con expresiones propias de la región. La
estudiante, centrada en escrituras, contratos y recibos, concluye que el relato es
confuso y que la persona ``no sabe explicar bien'' su problema. La pregunta que
guía la reflexión es cómo una forma de entender lo válido, lo confiable y lo bien
explicado puede hacer que otra manera de narrar un conflicto parezca insuficiente
en una situación jurídica.

% Técnica didáctica: foro académico (100 técnicas didácticas UnADM). Trazabilidad
% metodológica en comentario para no introducir metadiscurso en el texto visible.
\section{Diversidad cultural y lingüística como criterio de interpretación jurídica}

\subsection{El marco desde el que se juzga un relato}

El punto decisivo del caso no es que la persona explique mal su conflicto, sino
desde qué marco se juzga su relato. Cuando solo se reconoce como confiable lo
escrito, lo individual y lo expresado en lenguaje jurídico estándar, se corre el
riesgo de confundir la diferencia cultural con la falta de credibilidad. Como
advierte la comunicación intercultural, en toda interacción las personas no solo
transmiten datos, sino que narran desde sus marcos de autoridad, memoria y
pertenencia \citep{pech2016manual}. Para ordenar el análisis, tres preguntas guía
resultan pertinentes: ¿qué supuestos hacen que una narración comunitaria parezca
``confusa''?; ¿cómo distinguir entre una prueba insuficiente y una diferencia en la
forma de narrar?; y ¿qué debe hacer la persona profesional del Derecho para no
convertir la diversidad en desventaja jurídica?

\subsection{Sesgo cultural y distinción entre prueba y traducción}

El \textbf{sesgo cultural} puede identificarse en tres premisas y sus efectos profesionales. Considerar que solo lo documentado vale como prueba puede excluir indicios pertinentes; asumir que una narración lineal es más confiable que una narración colectiva puede producir una comprensión incompleta de los hechos; e interpretar las expresiones regionales como falta de claridad puede generar un trato potencialmente discriminatorio. En conjunto, estas premisas muestran cómo una lectura culturalmente no situada puede convertir la diferencia en una desventaja jurídica \citep{pech2016manual}.

Sobre la distinción entre prueba insuficiente y diferencia cultural, reconocer la
mediación cultural no significa idealizar la oralidad ni suponer que toda práctica
comunitaria sustituye los requisitos jurídicos. Un acuerdo de palabra, la presencia
de testigos comunitarios o la intervención familiar no vuelven verdadero un reclamo
por sí mismos, pero sí obligan a investigar antes de descartarlo. La perspectiva
intercultural aporta una corrección metodológica: interpretar los hechos
considerando la cosmovisión, la lengua y las instituciones comunitarias, en lugar
de imponer un único estándar de validez \citep{scjn2022protocolo}.

En cuanto a la actuación profesional, la respuesta no es descalificar el relato,
sino reconstruir el conflicto. La actuación profesional se concreta en una \textbf{entrevista jurídica intercultural}: primero, escuchar sin descalificar y permitir que la persona explique el conflicto desde sus propios referentes; después, contextualizar la lengua, las prácticas comunitarias y las formas locales de autoridad; luego, contrastar el relato con testigos y otros elementos de corroboración pertinentes; y, finalmente, argumentar la decisión distinguiendo entre una prueba insuficiente y una falta de traducción cultural. Esta secuencia permite incorporar el contexto sin sustituir la valoración jurídica de la prueba \citep{scjn2022protocolo}.
Por tanto, en un país cuya Constitución reconoce la composición pluricultural de la
nación y protege el derecho de las personas indígenas a ser asistidas por
intérpretes que conozcan su lengua y cultura, ignorar esa diversidad puede traducirse
en una barrera de acceso a la justicia \citep{cpeum2026,leyDerechosLinguisticos2024}.

\subsection{Intervención argumentada sobre el caso}

% ==============================================================================
% META (uso interno, no visible en el PDF): el producto solicitado es un FORO
% (técnica de las 100 Técnicas Didácticas de la UnADM, apertura -> participación
% -> coordinación -> cierre). El bloque siguiente reproduce, de forma textual e
% ÍNTEGRA, la intervención publicada (no un resumen). El texto que lo rodea NO
% nombra 'la actividad' ni 'el foro' como narrador externo: gravita sobre el tema,
% preparando la intervención (arriba) e interpretándola (abajo). El bloque usa
% tcolorbox con texto JUSTIFICADO, numeración de intervenciones a la izquierda y un
% botón de copia (adjunto .txt embebido) para extraer el texto.
% ==============================================================================

\begin{forobox}[title={Participación publicada en el foro}]
\hfill\foroCopyButton{foro-participacion-Actividad-2.txt}\\[-0.5em]
\phantomsection\label{foro-participacion}

\textbf{Asunto:} Diversidad cultural y lingüística en el derecho\par
\medskip
Hola a todas y todos: comparto mi participación al foro a partir del caso de la estudiante que atiende un conflicto por un terreno acordado de palabra ante la comunidad.
\medskip

\begin{enumerate}[leftmargin=1.6em,labelsep=0.6em,itemsep=0.5em,label=\textbf{\arabic*)}]
  \item \textbf{El problema no es que la persona ``explique mal'', sino desde qué marco se juzga su relato.} Cuando solo se acepta como confiable lo escrito, lo individual y lo dicho en lenguaje jurídico estándar, se confunde la diferencia cultural con la falta de credibilidad. La comunicación intercultural recuerda que las personas no solo transmiten datos: narran desde sus marcos de autoridad, memoria y pertenencia \citep{pech2016manual}. Por eso, lo que la estudiante lee como ``confusión'' puede ser una narración situada, comunitaria y no lineal.

  \item \textbf{Reconocer la mediación cultural no equivale a idealizar la oralidad.} No supone que un acuerdo de palabra baste como prueba plena; supone investigar antes de descartar: identificar la lengua y las prácticas comunitarias, ubicar a las personas testigos y buscar corroboración adecuada al contexto, distinguiendo entre una prueba insuficiente y una falta de traducción cultural. Como señala la Suprema Corte, juzgar con perspectiva intercultural implica ``reconocer la existencia de diversas formas de comprender el mundo y de resolver los conflictos, sin jerarquizarlas ni imponer una visión única de lo que es válido'' (Suprema Corte de Justicia de la Nación, 2022, p.~17). Ese criterio ofrece una vía para incorporar el contexto de personas, pueblos y comunidades indígenas en el análisis jurídico \citep{scjn2022protocolo}.

  \item \textbf{Lo vinculo con mi formación en Derecho.} Una entrevista o una valoración de hechos que no reconoce la mediación cultural y lingüística puede reproducir exclusión y vulnerar derechos lingüísticos reconocidos por la ley \citep{leyDerechosLinguisticos2024}. En un país cuya Constitución reconoce la composición pluricultural de la nación, ignorar esa diversidad puede traducirse en una barrera de acceso a la justicia \citep{cpeum2026}. Mi regla de actuación es preguntar antes de concluir, contextualizar antes de calificar y garantizar interpretación cuando haga falta.
\end{enumerate}
\medskip

\textbf{Pregunta para el grupo:} cuando atienden un relato que no encaja en el formato documental esperado, ¿qué hacen primero para no confundir la diferencia cultural con la falta de pruebas?\par
\medskip
Saludos y quedo atento a sus comentarios.
\medskip
\par\noindent\textbf{Referencias}\par
{\setlength{\parskip}{0.35em}%
\par\noindent\hangindent=1.6em\hangafter=1 Cámara de Diputados del H. Congreso de la Unión. (2024). \textit{Ley General de Derechos Lingüísticos de los Pueblos Indígenas}. https://www.diputados.gob.mx/LeyesBiblio/pdf/LGDLPI.pdf
\par\noindent\hangindent=1.6em\hangafter=1 Cámara de Diputados del H. Congreso de la Unión. (2026). \textit{Constitución Política de los Estados Unidos Mexicanos}. https://www.diputados.gob.mx/LeyesBiblio/pdf/CPEUM.pdf
\par\noindent\hangindent=1.6em\hangafter=1 Pech Salvador, C. E., Rizo García, M., \& Romeu Aldaya, V. (2016). \textit{Manual de comunicación intercultural: Una introducción a sus conceptos, teorías y aplicaciones} (2.\textsuperscript{a} ed.). Universidad Autónoma de la Ciudad de México.
\par\noindent\hangindent=1.6em\hangafter=1 Suprema Corte de Justicia de la Nación. (2022). \textit{Protocolo para juzgar con perspectiva intercultural: Personas, pueblos y comunidades indígenas}. Suprema Corte de Justicia de la Nación.
\par}
\end{forobox}

Leída en conjunto, la intervención muestra que la competencia intercultural no
debilita el rigor probatorio: lo ordena. La pregunta final apunta al punto que más
suele fallar en la práctica —el primer gesto de la escucha profesional—, porque de
ahí depende que la diversidad se lea como información jurídicamente relevante y no
como una deficiencia del relato.

\subsection{Diálogo entre pares sobre la valoración de la prueba}

Por privacidad académica se omite el nombre de la persona participante y se recupera
solo la idea pertinente. Esa participación sostiene, en síntesis, que la estudiante
del caso actúa bien al exigir documentos, pues el Derecho necesita pruebas, y que sin
papeles no había forma de avanzar. Frente a ese planteamiento, la réplica publicada en
el foro busca conservar el valor de la prueba sin desplazar el problema hacia la
persona usuaria.

\begin{forobox}[title={Retroalimentación publicada a otra participación}]
\phantomsection\label{foro-retroalimentacion}

Coincido en que la prueba documental es fundamental y conviene conservar ese punto.
Sin embargo, conviene distinguir dos operaciones profesionales que el caso tiende a fundir: por un lado, \textbf{exigir pruebas y corroboraciones pertinentes}, lo cual es plenamente legítimo; por otro, \textbf{concluir que la persona ``no sabe explicar''} por la forma de su relato, lo cual desplaza el problema hacia ella y no hacia el marco desde el que se le escucha. Por ello, un acuerdo de palabra ante la comunidad, con testigos y participación familiar, no es prueba plena, pero sí una pista de indagación que obliga a preguntar más antes de cerrar el caso. La perspectiva intercultural exige valorar ese contexto sin sustituir la corroboración jurídica \citep{scjn2022protocolo}. Como ejemplo, si la persona hablara una lengua indígena, la salida no sería descartarla por ``confusa'', sino garantizar interpretación conforme a los derechos lingüísticos \citep{leyDerechosLinguisticos2024}.\par
\medskip
\textbf{Pregunta para quien participó:} ¿cómo cambiaría tu conclusión si vieras la ausencia de papeles no como un límite de la persona usuaria, sino como una tarea de investigación que corresponde a quien ejerce la abogacía?
\medskip
\par\noindent\textbf{Referencias}\par
{\setlength{\parskip}{0.35em}%
\par\noindent\hangindent=1.6em\hangafter=1 Cámara de Diputados del H. Congreso de la Unión. (2024). \textit{Ley General de Derechos Lingüísticos de los Pueblos Indígenas}. https://www.diputados.gob.mx/LeyesBiblio/pdf/LGDLPI.pdf
\par\noindent\hangindent=1.6em\hangafter=1 Suprema Corte de Justicia de la Nación. (2022). \textit{Protocolo para juzgar con perspectiva intercultural: Personas, pueblos y comunidades indígenas}. Suprema Corte de Justicia de la Nación.
\par}
\end{forobox}

\clearpage
\section{Conclusión}

A mi juicio, la forma más sólida de resolver el caso es tratarlo como un problema de
sesgo jurídico y no de ``mala explicación''.

Los criterios de lo válido, lo confiable y lo bien explicado están culturalmente
situados; por ello, aplicarlos como un estándar único puede descalificar narrativas
orales, comunitarias y regionales jurídicamente relevantes. En consecuencia, sostengo
que la competencia intercultural no sustituye el rigor probatorio, sino que exige
investigar, contextualizar y valorar los hechos sin confundir la diferencia cultural
con falta de credibilidad \citep{scjn2022protocolo}. Distinguir entre falta de prueba
y falta de traducción cultural es, así, una condición de acceso efectivo a la
justicia cuando están comprometidos derechos lingüísticos y culturales
\citep{cpeum2026,leyDerechosLinguisticos2024}.

Por todo ello, considero que la postura profesional más consistente es investigar
antes de concluir, traducir sentidos antes de exigir tecnicismos y mantener abierta la
búsqueda de corroboración sin invalidar de inicio el relato de la
persona.\footnote{Para redactar y organizar este documento se utilizó una
herramienta de inteligencia artificial con el propósito de ordenar ideas y revisar
la redacción; su uso no sustituyó el análisis del caso, la selección y aplicación de
las fuentes ni la retroalimentación entre pares, que corresponden al autor.}

\clearpage
\nocite{bertely2011interaprendizajes,pech2016manual,scjn2022protocolo,diazAlonso2011construccion,cpeum2026,leyDerechosLinguisticos2024,oit1691989,unescoInterculturalCompetences2013}
\bibliography{interaprendizaje-en-ambientes-virtuales}

\end{document}
